\chapter{Ingegneria del Software}
Per definizione:
\begin{itemize}
    \item È una \textbf{disciplina} dell'ingegneria che riguarda tutti gli aspetti del software dai requisiti fino al mantenimento futuro.
    \item \textbf{CMU} dice che sia una forma di ingegneria che applica i principi dell'IT per ottenere una soluzione costi-efficace.
    \item \textbf{IEEE 610} invece la considera un approccio sistematico, disciplinato e quantificabile allo sviluppo e mantenimento.
\end{itemize}
Un progetto viene definito un successo che rispetta tre parametri di misurazione:
\begin{enumerate}
    \item La consegna è stata effettuata in \textbf{tempo}.
    \item Il progetto è rientrato nei \textbf{costi} preventivati.
    \item Ha prodotto un \textbf{output} che ha soddisfatto il cliente.
\end{enumerate}
Se uno di questi non è stato rispettato il progetto è \textbf{\textcolor{blue}{challenged}}, se nessuno si dice che ha \textbf{\textcolor{red}{fallito}}.
Tendenzialmente i progetti hanno successo quando coinvolgono l'utente, hanno richieste chiare e sono pianificati bene.
Altrimenti falliscono soprattutto per mancanza di input, chiarificazioni delle specifiche o cambiamento in corso d'opera:
\begin{enumerate}
    \item Buon \textbf{team di sviluppo}: buone skill, coordinamento, influenza, problem solving e confronto.
    \item Buona \textbf{pubblicità}: visione chiara del progetto, tempi decisionali ristretti, disponibilità.
    \item Buoni \textbf{spazi di sviluppo}: coinvolgimento dell'utente, ottimizzazione, esecuzione veloce.
\end{enumerate}

Gli \textbf{\textcolor{red}{errori}} sono generalmente generati dal codice, ma può capitare che si siano propagati da delle specifiche non chiare. Per ridurre tali \textbf{bug} si preferisce nei progetti ampi utilizzare parti di codice già sviluppati. \\

\[
\textit{Un ingegnere del software è una costruzione multipersonale di un software multi versione}.
\]

Eticamente per l'ingegneria del software ci sono \textbf{\textcolor{blue}{8} principi}:
\begin{enumerate}
    \item Pubblico interessato.
    \item Cliente e Datore.
    \item Prodotto.
    \item Giudizio sull'integrità integrità del prodotto.
    \item Gestione del processo etico.
    \item Professionalità.
    \item Colleghi.
    \item Pratica e continua formazione.
\end{enumerate}
Tali principi vengono utilizzati per la privacy, la proprietà
intellettuale, il riuso del codice, la risoluzione dei problemi.

\vspace*{0.02\textheight}
\textbf{Alan Davis} nel \textbf{1994} ha definito alcune delle regole per sviluppare un software professionale:
\begin{enumerate}
    \item Qualità al primo posto, dato che è possibile fare software di qualità.
    \item Fornire il prodotto al cliente il prima possibile.
    \item Inquadrare il problema prima di iniziare a segnare le specifiche.
    \item Valutare eventuali alternative di progettazione.
    \item Usare un appropriato modello di processo.
    \item Usare linguaggi differenti per diverse fasi dello sviluppo.
    \item Mettere le tecniche davanti agli strumenti per applicarle.
    \item Preferire fare le scelte giuste, rispetto che quelle più veloci.
    \item Una buona gestione è più importante di una buona tecnologia.
    \item Le persone sono la chiave del successo, quindi prendersi responsabilità e seguire il progetto con cura.
\end{enumerate}

\vspace*{0.02\textheight}
\textbf{Anthony Wasserman} nel \textbf{1996} ha proposto alcuni concetti base per tale disciplina:
\begin{enumerate}
    \item \textbf{Astrazione:} una tecnica utilizzata per la progettazione che permette di generalizzare concetti.
    \item \textbf{Analisi e metodi/notazioni di sviluppo:} standardizzare i requisiti e le notazioni per comunicare universalmente.
    \item \textbf{Prototipare la UI:} permette di inquadrare le preferenze dell'utente prima di iniziare lo sviluppo.
    \item \textbf{Modularità e architettura:} permette di aumentare la qualità e la manutenibilità del software.
    \item \textbf{Riutilizzo:} è una pratica nell'industria software per permettere di allungare la vita del progetto.
    \item \textbf{Ciclo di vita:} deve essere considerato dall'inizio alla fine.
    \item \textbf{Metriche di processo:} permette di verificare spesso l'andamento del software.
    \item \textbf{Strumenti e ambienti integrati:} essenziali per sviluppare software complessi.
\end{enumerate}

\section{Processi Software}
Per sviluppare un buon software c'è bisogno di varie \textbf{task} svolte da diverse persone e coordinate da un terzo:
\begin{itemize}
    \item \textbf{\textcolor{brown}{CODE N FIX:}} modello più semplice di programmare da solo che una volta soddisfatte le specifiche si fa test e debug.
    \item \textbf{\textcolor{brown}{WATERFALL:}} modello a cascata che segue gli step principali a cui si aggiunge il mantenimento del software:
    \begin{itemize}
        \item L'\textbf{output} del task precedente è l'\textbf{input} del successivo ed è utile per descrivere lo stato attuale del progetto.
        \item Purtroppo l'interazione con l'utente avviene solo nella fase requisiti ed è \textbf{document-driven} perché ne fa molti.
    \end{itemize}
    \item \textbf{\textcolor{brown}{INCREMENTALE:}} ogni feature viene sviluppata separatamente poi unita tramite integrazione e testata:
    \begin{itemize}
        \item I vantaggi sono sviluppare per primi i componenti più importanti e contenere il rischio se c'è un problema.
        \item Può essere utilizzato per sviluppo \textbf{multi-componente}, oppure \textbf{multi-release} tramite varie versioni del prodotto.
        \item In pratica utilizza il metodo \textbf{Divide and Conquer} ed è precursore della tecnica attuale \textbf{\textcolor{blue}{CI/CD}}.
    \end{itemize}
\end{itemize}

